
\documentclass[conference]{IEEEtran}
\IEEEoverridecommandlockouts
\usepackage{cite}
\usepackage{amsmath,amssymb,amsfonts}
\usepackage{algorithmic}
\usepackage{graphicx}
\usepackage{textcomp}
\usepackage{xcolor}
\usepackage{booktabs}
\usepackage[hidelinks]{hyperref}
\def\BibTeX{{\rm B\kern-.05em{\sc i\kern-.025em b}\kern-.08em
    T\kern-.1667em\lower.7ex\hbox{E}\kern-.125emX}}

\newcommand{\tofill}[1]{{\color{red}\textbf{[FILL: #1]}}}

\begin{document}

\title{Grading Prostate Cancer Severity from Whole-Slide Biopsies on the PANDA Challenge}

\author{
\IEEEauthorblockN{Shashanth Markabundu}
\IEEEauthorblockA{
\textit{AI in Medical Imaging} \\
\textit{Radboud University} \\
Nijmegen, Netherlands \\
shashanth.markabundu@ru.nl
}
\and
\IEEEauthorblockN{Francisca  Carneiro}
\IEEEauthorblockA{
\textit{AI in Medical Imaging} \\
\textit{Radboud University} \\
Nijmegen, Netherlands \\
francisca.carneiro@ru.nl
}
\and
\IEEEauthorblockN{Goja Modic Sila}
\IEEEauthorblockA{
\textit{AI in Medical Imaging} \\
\textit{Radboud University} \\
Nijmegen, Netherlands \\
goja.modicsila@ru.nl
}
}

\maketitle

\begin{abstract}
Prostate cancer is among the most common cancers in men worldwide, and its
management depends on the ISUP grade (0 to 5) assigned to biopsy whole-slide
images via the Gleason system. Manual grading is decisive but suffers from
substantial inter-observer disagreement, which makes consistent automation
valuable. The PANDA challenge poses two characteristic difficulties. Each image
is a gigapixel whole-slide scan that must be sampled efficiently, and the
training labels are imperfect because, unlike the multi-rater test set, they were
graded by a single pathologist. We compare two modelling paradigms on an
identical tiling pipeline and an identical ordinal target evaluated with
quadratic weighted kappa (QWK). The first is the convolutional recipe of the 2020
challenge era, an EfficientNet-B0 trained on tissue-tile montages. The second is
the 2024 foundation-model recipe, a frozen pathology vision transformer (UNI)
used as a tile encoder with a small gated-attention multiple-instance-learning
head. On a held-out fold of 2124 slides the CNN reached QWK 0.876 and the
foundation model 0.859, both weaker on Radboud than on Karolinska. A
deterministic Gleason to ISUP audit of all 10616 training slides found a single
inconsistent label, locating the residual noise in grading subjectivity rather
than data entry and explaining why errors cluster on adjacent grades. Critically,
when the foundation model is trained on one centre and tested on the other its QWK
falls to 0.485 to 0.649, showing that strong in-distribution accuracy can mask poor
cross-institutional generalisation. We also document the offline, notebooks-only
inference engineering that was the project's largest practical hurdle.
\end{abstract}

\begin{IEEEkeywords}
computational pathology, prostate cancer, Gleason grading, foundation models,
multiple instance learning, domain shift, quadratic weighted kappa
\end{IEEEkeywords}

\section{Introduction}

Prostate cancer grading on biopsy whole-slide images (WSIs) is summarised by the
ISUP grade, an ordinal severity score from 0 (benign) to 5 (most aggressive)
derived from the Gleason pattern. Grading is clinically decisive but suffers from
substantial inter-observer disagreement among pathologists, which makes it an
attractive target for automation and, simultaneously, a hard one to supervise
the ``ground truth'' is itself uncertain. The PANDA challenge \cite{bulten}
released the largest public corpus for this task and established quadratic
weighted kappa (QWK) as the metric, because QWK rewards predictions that are close
on the ordinal scale and penalises being several grades off far more than being
one grade off.

Two paradigms bracket this problem. The first, behind the strongest 2020 PANDA
solutions, trains a convolutional network from scratch on tissue tiles, typically
stitched into a montage so one forward pass sees a representative sample of the
slide \cite{tan,baselines}. Its strength is end to end features learned for exactly
this task. Its weaknesses are documented such that models are data-hungry and
demonstrably encode site-specific signatures (Howard \textit{et al.} showed CNNs
can predict the submitting institution from histology alone \cite{howard}), and
stay sensitive to stain and scanner variation unless explicitly normalised
\cite{tellez}. The second paradigm, dominant since 2024, uses a foundation
model. A vision transformer self-supervised on $10^5$+ slides, frozen as a generic
tile encoder with a light trained head; UNI \cite{chen} and Virchow \cite{virchow}
are representative. Its strength is reusable representations from little labelled
data, and its hoped for advantage is robustness from having seen far more
morphology than any single challenge set. Its weakness is that the frozen encoder
never adapts to the target domain, and a recent re-annotation benchmark warns such
models reach high in-distribution accuracy by keying on specimen-specific artefacts
while transferring poorly across slides \cite{pandaplus}. We adopt UNI as the
exemplar openly accessible, a ViT-L backbone trained on a large multi-institution
corpus, with strong reported pathology benchmarks.

The PANDA training data is drawn from two institutions, Radboud and Karolinska
\cite{strom}, which differ in scanner, pathologist panel and even mask annotation
scheme. This built-in domain shift turns ``which paradigm generalises better
across centres?'' into an answerable question rather than a solved one. We
therefore make three contributions: (i) a data inspection that quantifies class
imbalance, the two-centre distribution gap, label-scheme differences and a direct
Gleason-to-ISUP label consistency audit (ii) a controlled comparison of the CNN
and the frozen foundation model plus MIL head on an identical pipeline, ordinal
target and metric and (iii) a cross-centre robustness stress-test of the
foundation model, with a clinically framed error analysis, showing that strong
in distribution QWK masks a large generalisation gap across institutions. We
additionally document the offline-inference engineering that the challenge's
notebooks only submission format demanded.

\section{Data and Its Challenges}

All experiments use the PANDA training set: 10{,}616 prostate-biopsy WSIs with an
ISUP grade per slide and, for 10{,}516 of them, a pixel-level mask. Inspecting the
data before modelling motivates every later design choice.

\textbf{Class imbalance and ordinality.} The ISUP grades are imbalanced toward
benign and low-grade slides. Plain accuracy would reward predicting the majority
grade. This is the first reason the metric is QWK and the second reason we adopt
an ordinal target (Section~III) rather than treating the six grades as unordered
classes predicting 4 when the truth is 5 must cost less than predicting 0.

\textbf{Two centres are a built-in domain shift.} The slides come from Radboud
(5{,}160) and Karolinska (5{,}456), on different scanners and pathologists, and
the two centres have markedly different grade distributions (Fig.~\ref{fig:dist}):
Karolinska is dominated by benign and ISUP~1 (about 35\% and 33\% of its slides),
whereas Radboud is far flatter and rises again at ISUP~5 (about 19\%). A model can
thus exploit centre-specific colour or the centre prior as a shortcut for
biology---the basis of our cross-centre experiment.

\textbf{Label-scheme asymmetry.} Where masks exist (10{,}516 slides), the centres
annotate incompatibly: Radboud labels individual glands with five tissue classes
(stroma, benign epithelium, Gleason 3/4/5) while Karolinska distinguishes only
benign versus cancer---so any mask-based analysis must treat the schemes separately.

\textbf{Train/test artefacts and tissue content.} Training slides carry stray pen
marks absent from the test set---a train/test difference a model could key on---and
tissue fraction varies widely (some biopsies are over 98\% background), making the
most-tissue tile-selection step essential.

\textbf{Label-consistency audit.} ISUP is derived \emph{deterministically} from the
Gleason score by fixed rules ($3{+}3\!\rightarrow\!1$, $3{+}4\!\rightarrow\!2$,
$4{+}3\!\rightarrow\!3$, and so on). Computing the expected ISUP for every slide and
comparing it with the recorded ISUP across all 10{,}616 slides yields only
\textbf{one} inconsistency (0.01\%; a Karolinska $4{+}3$ slide recorded as ISUP~2
instead of 3). The recorded labels are therefore internally clean: the
well-documented PANDA noise comes from \emph{grading subjectivity}, not data-entry
error. This is compounded by a train/test asymmetry the organisers note---the test
slides were graded by multiple pathologists whereas the training labels rely on a
single rater---so the training target is inherently noisier than the benchmark it is
scored against. Because the labels are internally consistent, we expect, and later
observe, errors concentrated on \emph{adjacent} grades where pathologists themselves
disagree.

\begin{figure}[tb]
\centerline{\includegraphics[width=\columnwidth]{fig_distribution.png}}
\caption{Slides per centre (left) and ISUP distribution within each centre
(right). The two centres differ sharply---Karolinska skews benign/low-grade,
Radboud is flatter and higher-grade-heavy---a shortcut a model can exploit.
Produced by the EDA notebook.}
\label{fig:dist}
\end{figure}

\section{Methods}

\subsection{Experimental overview}
The data inspection drives the design: imbalance and ordinality motivate QWK and an
ordinal target, and the two-centre shift motivates the cross-centre study. All
experiments share a \emph{single pipeline}---a from-scratch CNN baseline
(Approach~A), a frozen-foundation-model arm on the same tiles and target
(Approach~B), and a cross-centre stress-test---so any difference between models is
attributable to the encoder, not the preprocessing.

\subsection{Shared tiling pipeline}
Each WSI is read at openslide level~1; the tissue is split into $256\times256$
tiles and the 36 tiles with the most tissue are kept. This expensive step is run
\emph{once} and cached: for the CNN the 36 tiles are stitched into a single
$6\times6$ ($1536\times1536$) montage JPEG; for the foundation model each slide is
stored as a $36\times1024$ embedding array. Caching converts the dominant
cost---decoding multi-level TIFFs every epoch---into a one-time operation, which is
what makes repeated training feasible within the compute budget. One subtlety
matters: the CNN montages are colour-inverted ($255-\text{pixel}$, following the
public baseline) whereas the foundation model expects ordinary H\&E,
ImageNet-normalised and resized to $224$, \emph{not} inverted. Mixing these
conventions silently destroys accuracy, so they are kept strictly separate.

\subsection{Ordinal target and metric}
For both models the grade $k$ is encoded as a cumulative vector of length
five---the first $k$ entries set to one---and the network emits five independent
logits trained with binary cross-entropy. At inference the predicted grade is the
rounded sum of the five output sigmoids. This cumulative formulation respects grade
ordering and directly targets QWK, which is why it is preferred over a six-way
softmax; the public baselines report it lifting the leaderboard from $\approx0.79$
(softmax) to $\approx0.87$ (ordinal BCE)---a deliberate, literature-grounded choice
rather than an inherited default.

\subsection{Approach A: from-scratch EfficientNet-B0}
The baseline is an EfficientNet-B0 \cite{tan} with a five-output linear head,
initialised from ImageNet weights and fine-tuned end-to-end on the inverted
montages. Training uses Adam (initial learning rate $3\times10^{-4}$), one warm-up
epoch then cosine annealing, mixed precision, batch size~4, and orientation-only
augmentation (transpose and flips), appropriate because tissue orientation carries
no diagnostic meaning; two-view test-time augmentation is applied at inference. The
original public recipe trains for 30 epochs, but on a single Kaggle GPU this took
over 11 hours and had to be cut to 20 epochs to fit the session and weekly quota.
The validation QWK was still rising slowly at that cut-off, so the reported 0.876
is a floor rather than a converged ceiling---a longer schedule would very likely
score higher.

\subsection{Approach B: frozen UNI encoder + gated-attention MIL}
The second model uses UNI \cite{chen}, a ViT-L pathology foundation model, frozen
as a tile encoder that maps each of the 36 tiles to a 1024-dimensional embedding.
A slide is thus a bag of 36 instances, and we train only a gated-attention MIL head
\cite{ilse}: a shared projection, a gated attention mechanism that scores each
tile, a softmax over tiles, and an attention-weighted sum fed to the same
five-output ordinal head as Approach~A. Because only the head is trained and the
encoder is frozen, it converges in minutes from the precomputed embeddings, and the
attention weights are directly interpretable as which tiles drove a prediction.

\subsection{Learning strategy and overfitting control}
We monitor the learning process directly: BCE loss and validation QWK (overall and
per centre) are logged every epoch. Three observations guide the strategy. (i)
Training and validation loss fall together with only a small, stable gap, so the
CNN is \emph{not} overfitting---ImageNet initialisation, orientation augmentation
and best-checkpoint selection (the highest-QWK epoch, not the last) all regularise.
(ii) Validation QWK rises steeply and then flattens, but at the 20-epoch compute
cut-off (Section~III-D) it had \emph{not} fully plateaued; relative to the 30-epoch
public recipe the model is mildly \emph{underfit}, which is the clearest single
lever for a higher score. (iii) The per-centre QWK curves separate early, an
in-training signal of the domain gap quantified later. The MIL head needs different
control: few trainable parameters on a frozen encoder cannot overfit 8{,}500 slides,
so dropout ($p=0.25$) suffices. Validation uses fold~0 of a stratified 5-fold split
(\texttt{random\_state}=42), 2{,}124 held-out slides; the cross-centre study
re-trains on one centre and tests on the whole other centre, both directions.

\subsection{Inference under the challenge constraints}
The single largest practical obstacle was \emph{inference}. PANDA is a
notebooks-only competition: submissions run as Kaggle notebooks with the internet
disabled and the hidden test set mounted only at run time, and our initial
experiments would not run there for several reasons, each needing a fix. The 2020
baseline's dependencies no longer matched the current Kaggle image (its EfficientNet
source package and multi-level-TIFF reader both failed), so the pipeline was
re-based on the preinstalled \texttt{timm} backbone and \texttt{openslide} reader.
Three submission constraints then bind. First, no package may be installed offline,
so inference uses only those libraries and loads weights from attached Datasets.
Second, the hidden test set is \emph{not} pre-tiled, so inference must read and tile
each test TIFF on the fly---the very decoding cost the cache avoids---with tiling and
inversion identical to training, or the weights misbehave. Third, UNI's gated weights
cannot be fetched offline, so they were downloaded once and shipped as a private
Dataset. Reconciling these---and fixing a train/inference weight-name mismatch and a
tiler bug---is what produced valid offline submissions for both arms.

\section{Results}

\subsection{Training behaviour}
The CNN's validation QWK climbed from 0.77 at epoch~1 to 0.79 by epoch~3 and on to
the reported 0.876, with training and validation loss falling together (no
overfitting signature) and the best-QWK epoch retained. As anticipated in
Section~III-D, the curve had not plateaued at the 20-epoch compute cut-off, so
0.876 is a floor rather than a ceiling: the 30-epoch public recipe this arm follows
scores higher, and a full schedule would recover part of that gap. The per-centre
split was already visible in the first epochs (epoch~3: Karolinska 0.82 versus
Radboud 0.68), foreshadowing the domain gap quantified below. The MIL head, by
contrast, converges in minutes on the precomputed embeddings.

\subsection{In-centre comparison}
On the held-out fold the from-scratch CNN reaches QWK 0.876 and the frozen
foundation model 0.859; on the hidden leaderboard they score 0.898/0.875 and
0.868/0.845 (private/public), preserving the validation ordering. A separate
submission using the public pre-trained weights scored 0.915, but we report only
our own model because the challenge requires training to be reproducible; the gap
to 0.915 is consistent with our shorter, single-fold, compute-limited schedule
(Section~III-D). The CNN also makes fewer clinically serious errors---43 missed
cancers and 72 false alarms, versus 62 and 120 for the foundation model---where a
\emph{missed cancer} is a slide predicted benign that is truly cancerous and a
\emph{false alarm} the reverse.

\subsection{Per-centre performance}
Evaluated separately on each centre ($n{=}1{,}111$ Karolinska, $1{,}013$ Radboud),
both jointly trained models are stronger on Karolinska than on Radboud:
EfficientNet-B0 scores 0.891 versus 0.835 and the foundation model 0.849 versus
0.826. Radboud is the harder centre for both, consistent with its finer five-class
annotation scheme and greater grade diversity.

\subsection{Cross-centre robustness of the foundation model}
Retraining the foundation-model pipeline on one centre and testing on the entire
other centre, in both directions (test sets of $n{=}5{,}160$ Radboud and
$n{=}5{,}456$ Karolinska slides). Performance collapses relative to the
in-distribution 0.85 range: QWK falls to 0.649 when training on Karolinska and
testing on Radboud, and to 0.485 in the reverse direction. The transfer is also
asymmetric---a Karolinska-trained model generalises to Radboud markedly better than
the reverse---mirroring the per-centre difficulty ordering, and the cross-centre
confusion matrices show the model mis-binning grades toward the source centre's
prior rather than scattering randomly.

\subsection{Error analysis}
The validation confusion matrices (Fig.~\ref{fig:cm}) show both models' errors
lying overwhelmingly on the diagonal's immediate neighbours: for the foundation
model 92.7\% of errors are within a single grade, and for the CNN none of the 30
largest errors coincide with the single audited label inconsistency. Because the
audit found only one inconsistent label in 10{,}616 slides, these errors cannot be
data-entry noise; they fall on genuinely ambiguous, adjacent-grade boundaries where
pathologists disagree---the expected behaviour of a well-trained model against a
subjective reference standard. This interpretation directly shapes the next
experiments: because the error is adjacent-grade rather than catastrophic, the
productive levers are ordinal-target label smoothing and a longer training schedule
(Section~III-D), not a change of architecture.

\begin{figure}[tb]
\centerline{\includegraphics[width=0.84\columnwidth]{fig_confusion.png}}
\caption{Validation confusion matrices for the two models. Errors concentrate on
adjacent grades, consistent with inter-observer grading variability rather than
label corruption. Produced by the EDA notebook.}
\label{fig:cm}
\end{figure}

\section{Discussion}

\textbf{The cross-centre collapse is the main finding.} In-centre, the CNN modestly
beat the frozen foundation model on QWK and serious errors---credible, since the CNN
is fine-tuned end-to-end while UNI is frozen and never adapts to PANDA. The more
important result is that UNI, despite its $\approx$0.85 in-distribution QWK, drops to
0.485--0.649 across centres: high in-distribution accuracy masks poor
cross-institutional generalisation, the frozen encoder still keying on
centre-specific appearance rather than transferable biology.

\textbf{Reasons for failure.} \emph{Label subjectivity}: the audit proves labels
are internally consistent, so adjacent-grade errors reflect real pathologist
disagreement---a ceiling no model crosses without multi-rater labels.
\emph{Domain shift}: the centres differ in stain and scanner, so the cross-centre
drop reflects features that do not transfer. \emph{Information loss}: level-1 tiling
and JPEG compression of the montage discard detail fine Gleason discrimination may
need.

\textbf{Broader context.} The original PANDA paper flagged that
development-to-external distribution shift caused over-diagnosis \cite{bulten}, and
recent re-annotation benchmarks report the same artefact-keying \cite{pandaplus}.
Our 0.485--0.649 cross-centre QWK quantifies that effect inside the public training
data.

\textbf{Limitations and a results-driven plan.} Two gains need no retraining, since
the arms share an ordinal head: a weighted average of their continuous scores, and
QWK-optimised rounding thresholds fitted on the validation fold. The deeper
questions need new runs, in a findings-driven sequence. First, because the
foundation model collapsed cross-centre, the immediate experiment is the
\emph{identical} cross-centre protocol on the CNN---testing whether task-specific
training confers the robustness the frozen encoder lacked, the single most
informative missing comparison. Second, Macenko normalisation on both arms would
isolate how much of the gap is colour rather than morphology. Third, a full-length
schedule would quantify the headroom the 20-epoch limit cost. Remaining
limitations---single-fold evaluation, a compressed montage, and a frozen rather than
fine-tuned encoder---bound the present comparison.

\section{Conclusion}
On PANDA ISUP grading with an identical pipeline and ordinal target, a from-scratch
EfficientNet-B0 (QWK 0.876, compute-limited to 20 epochs) modestly outperformed a
frozen UNI encoder with a gated-attention MIL head (0.859) in-centre. A deterministic
label audit located PANDA's residual noise in grading subjectivity, explaining both
models' adjacent-grade errors. Stress-testing UNI across centres exposed a large
generalisation gap (0.485--0.649 versus $\approx$0.85 in-distribution): strong
in-centre accuracy does not imply cross-institutional robustness. Meeting the
offline, notebooks-only inference format was the largest engineering hurdle.

\section*{Code Availability}
The repository
contains eight notebooks in run order. Tile caching, EDA and label audit,
EfficientNet training, a validation-prediction dump, EfficientNet inference, UNI
embedding extraction, MIL-head training with the cross-centre study, and UNI
inference plus a README documenting the Kaggle-Dataset dependencies and the exact
run settings (GPU, internet, attachments) for each notebook. Both inference
notebooks reproduce \texttt{submission.csv} offline from saved weights. Code repository:
\href{https://github.com/Shashaboii/AIMI_Panda_Challenge}{GitHub Repository}

\begin{thebibliography}{00}
\bibitem{bulten} W. Bulten, K. Kartasalo, P.-H. C. Chen, \textit{et al.},
``Artificial intelligence for diagnosis and Gleason grading of prostate cancer: the
PANDA challenge,'' \textit{Nature Medicine}, vol.~28, pp.~154--163, 2022.
\bibitem{chen} R. J. Chen, T. Ding, M. Y. Lu, \textit{et al.}, ``Towards a
general-purpose foundation model for computational pathology,'' \textit{Nature
Medicine}, vol.~30, pp.~850--862, 2024.
\bibitem{howard} F. M. Howard, J. Dolezal, S. Kochanny, \textit{et al.}, ``The
impact of site-specific digital histology signatures on deep learning model
accuracy and bias,'' \textit{Nature Communications}, vol.~12, art.~4423, 2021.
\bibitem{tellez} D. Tellez, G. Litjens, P. B\'andi, \textit{et al.}, ``Quantifying
the effects of data augmentation and stain color normalization in convolutional
neural networks for computational pathology,'' \textit{Medical Image Analysis},
vol.~58, art.~101544, 2019.
\bibitem{virchow} E. Vorontsov, A. Bozkurt, A. Casson, \textit{et al.}, ``A
foundation model for clinical-grade computational pathology and rare cancer
detection,'' \textit{Nature Medicine}, vol.~30, pp.~2924--2935, 2024.
\bibitem{ilse} M. Ilse, J. M. Tomczak, and M. Welling, ``Attention-based deep
multiple instance learning,'' in \textit{Proc. ICML}, 2018, pp.~2127--2136.
\bibitem{tan} M. Tan and Q. V. Le, ``EfficientNet: rethinking model scaling for
convolutional neural networks,'' in \textit{Proc. ICML}, 2019, pp.~6105--6114.
\bibitem{strom} P. Str\"om, K. Kartasalo, H. Olsson, \textit{et al.},
``Artificial intelligence for diagnosis and grading of prostate cancer in biopsies:
a population-based, diagnostic study,'' \textit{Lancet Oncology}, vol.~21, no.~2,
pp.~222--232, 2020.

\end{thebibliography}

\end{document}
